\textbf{Время:} резерв

\textbf{Тезис:} параметр λ получен вместе с радиусами и зависит от выбранной выборки пар.

\textbf{На экране:} слева λ по kT, справа λ по быстроте, единая цветовая легенда внизу.

\textbf{Текст:} Это параметр лямбда текущих гауссовых описаний корреляции. Слева отборы по среднему поперечному импульсу пары, справа — по её быстроте в системе центра масс. Цвета имеют тот же смысл центральности, что у радиусов.

Параметр задаёт силу обменного члена в принятом описании. Он не заменяет пространственный размер: размер определялся по ширине структуры. Зависимость лямбда от отбора нужно рассматривать вместе с качеством описания и особенностями источника.

\textbf{Справка:} нижние правые панели kt\_c\_all\_graphs\_neg.jpg и c\_all\_graphs\_neg.jpg. Предоставленные изображения сохранены без изменения чисел. Значение λ нельзя само по себе объявлять измерением когерентности. Для проекций трёхмерной корреляции высота видимого пика зависит также от усреднения остальных компонент q; равенство λ=C(0)−1 относится к полной принятой гауссовой форме без дополнительных множителей.

\textbf{Источники:} два текущих JPEG, авторские подтверждения и манифесты; общие ограничения интерпретации λ — Lisa et al. (2005), §2.6.
